\section{Open questions}

\begin{enumerate}

\item \textbf{CDR quantitative reproduction on the Fig~5 benchmark.}
\emph{Basis:} This replication verified ZNE and PEC quantitatively but only checked CDR is present in the package. CDR is one of Mitiq's three headline primitives and its comparative accuracy/cost trade-off on the exact benchmark used for PEC (Fig~5) is unstated in the paper and untested here.
\emph{Next steps:} Add \texttt{work/rep\_cdr.py} using \texttt{mitiq.cdr.execute\_with\_cdr} on the same $H;X;CNOT$ circuit and \texttt{DensityMatrixSimulator} executor at $p=0.1$. Sweep \texttt{num\_training\_circuits} in $\{10, 50, 200\}$ and \texttt{fit\_function} in \{linear, exponential\}. Report mean $|err|$ over 10 seeds side-by-side with the existing PEC and ZNE numbers, plus wall-clock and total simulator calls, to yield a three-way cost/accuracy triangle.

\item \textbf{Multi-backend quantitative agreement (Cirq vs Qiskit-Aer).}
\emph{Basis:} The paper advertises Cirq/Qiskit/pyQuil/Braket support (C5). This replication only exercised the Cirq path; the Qiskit path is a separate code route through Mitiq's converters and could regress silently without affecting Cirq numbers.
\emph{Next steps:} Install \texttt{qiskit} + \texttt{qiskit-aer} in a parallel venv. Rebuild the Fig~5 circuit in Qiskit, wrap a \texttt{NoiseModel} with per-gate \texttt{depolarizing\_error(0.1)}. Feed it through \texttt{mitiq.pec.execute\_with\_pec} and \texttt{mitiq.zne.execute\_with\_zne}. Assert that the resulting mean $|err|$ matches the Cirq numbers to within Monte-Carlo variance.

\item \textbf{Mitiq ZNE vs newer mitigation techniques.}
\emph{Basis:} The paper is from 2020. Virtual distillation (Huggins 2021) and symmetry-based verification have since become standard comparisons. Whether Mitiq's ZNE remains competitive on canonical benchmarks or the newer methods dominate at the same sampling cost is an open, actionable question.
\emph{Next steps:} Extend \texttt{work/rep\_zne.py}: for each of the 20 RB circuits, additionally compute (a) a two-copy virtual-distillation estimator implemented from scratch in Cirq, (b) a symmetry-verification post-selection using the Z-parity conservation of the RB Clifford group. Report mean $|err|$, sample count, and classical post-processing time for each method vs unmitigated and vs ZNE.

\item \textbf{Real-hardware runtime overhead on IBM Quantum.}
\emph{Basis:} C5 (hardware runs) was not exercised. The paper reports hardware plots (Fig~3) but does not tabulate wall-clock or shot-count overhead in a form that lets a user predict cost on IBM's current queue-based system.
\emph{Next steps:} Using IBM Quantum's free-tier open-plan devices, run the Fig~5 circuit with and without ZNE (scale factors 1, 3, 5) at 8192 shots. Log job queue-time, circuit count, and calibration age. Compare per-mitigated-estimate wall-clock to the simulator baseline collected here.

\item \textbf{ML-in-the-loop mitigation-parameter selection.}
\emph{Basis:} Mitiq exposes a calibrator API but the paper uses fixed defaults (Richardson extrapolation, random local folding). Whether per-circuit adaptive parameter selection --- driven by a small classical model trained on circuit descriptors --- yields a real accuracy win is untested here and unresolved in the current mitigation literature.
\emph{Next steps:} Build a corpus of $\sim 200$ random 2--4 qubit circuits with varying depth and 2-qubit-gate density. For each, evaluate ZNE with a small grid of \texttt{(scale\_factors, extrap\_method)} combinations and record the best combo. Train a scikit-learn regressor mapping (\texttt{num\_qubits}, \texttt{depth}, \texttt{2q\_gate\_fraction}) $\to$ best combo. On a held-out set, compare the adaptive-choice mean $|err|$ against Mitiq's default. If it wins by $>1.5\times$ on average, propose upstreaming a \texttt{mitiq.calibrator.LearnedSelector}.

\end{enumerate}
